\section{An explicit cutoff at size 10,000}
\label{sec:paired-fourier-cutoff}

\textbf{Theorem.} For every \(n\ge10{,}000\) and \(3\le k\le n\), the
nonzero roots of \(E_{n,k}\) and \(E_{n+1,k}\) are simple, negative and
strictly interlace, with the next row supplying both outer nonzero roots.
Zero is a simple common root.

\paragraph{Paired residues for \(4\le n/k\le48\).}
Use the notation of Section~\ref{sec:degree-block-cutoff} and assume
\(nS^2\ge6/5\). For \(kS\ge10\) the joint value-and-derivative argument
applies with \(k\ge208\), using exponent 48 in the Bernoulli bound for
the half-powered off-diagonal term and 96 in the ordinary kernel.
For \(kS\le10\), retain the whole principal residue multiplier
\[
 M_1=d_\infty Q_1(z_1^{-1})
     =z_1(S/\pi)e^{-H(s+i\pi)},\qquad
 d_\infty=\exp\!\left(-\sum_{j=1}^{k-1}\frac1j\right),\quad
 s=\pi\cot\theta ,
\]
where \(H(w)=\sum_{b\ge1}b^{-1}e^{-bw}\sum_{j=0}^{k-1}(bw)^j/j!\).
In each bin of complementary degree \(T\), subtracting that
continuation replaces \(d_{n-T}\) by \(d_{n-T}-d_\infty\),
with the coefficient of \(D-1\) interpreted as zero below degree
\(k\). The same subtraction applies in the insertion's residual row.
The difference decays factorially in the number of full degree blocks.
Absolute convergence, including that of the angular derivative, allows
the tail after bin 240 to be bounded by evaluating the nonnegative
majorants at \(9/8\). The differentiated term uses the signed
frequency \(n-T\) and the separate multiplier derivative \(M_1'\).

The gamma-density bound on the vertical pole segment gives
\(\pi|H'|<0.18\). A harmonic sum and a Poisson lower-tail estimate
give \(2k|M_1|>0.47\), and its phase has cosine greater than \(0.97\).
The rational replay bounds both insertion vectors within
\(0.97(2k|M_1|)\) of
\(2k(\Re(M_1 z_1^n),\Re(-iM_1 z_1^n))\). Projection onto the
unrotated phase vector gives the integer sample signs and the
consecutive signs at old roots. The exterior ranks, reciprocity and
degree count then give strict interlacing in this range.
\href{https://github.com/ProgramComputer/open-math-research/tree/main/topics/long-cycle-excedances/paired-residues}{Proof and exact replay}

\paragraph{A Fourier contrast for \(n/k\ge44\).}
On a common pole-free circle \(w=Re^{i\varphi}\), put
\(f=e^G/(1+e^{w-s})\), where \(G\) is the truncated logarithm from
the common-contour proof. Write \(f_j\) for its \(j\)-th angular Fourier
coefficient, \(w_p=s+ip\pi\), \(r=|w_1|\) and
\(a_p=R^k/w_p^{k-1}\). The exact Taylor remainder gives
\[
 \partial_\varphi f
 =i e^{ik\varphi} f \sum_p\frac{a_p}{w_p-w}.
\]
Two integrations by parts give the factor \(1/(j(j-k))\) in
the circle error. For the consecutive contrast \((r/R) f_{n+1}-C f_n\),
first split each kernel as
\[
 \frac{r/w-C}{w_p-w}
   =\frac{z_p-C}{w_p-w}+\frac{z_p}{w},\qquad z_p=r/w_p .
\]
The first term retains the factor \(S\) through
\(|z_p-C|\le |p|S\). The second gives a lower Fourier coefficient
and is bounded by the same two integrations.

For \(k\ge100\), the radial separation and the contraction of the
Taylor numerator bound the normalized errors in terms of \(x=n/k\),
uniformly in the row size. The factor \(n(n-k)\) in the integrations
absorbs the two factors of \(k\) from the contour and kernel estimates.
The thresholds \(4\le k<100\) use direct common-circle bounds.
The earlier all-size theorem covers \(k=3\). For \(k\ge4\),
combining the circle and enclosed-pole errors gives
\[
 |Y_n-\cos\Psi|<0.20,\qquad
 |Y_{n+1}-CY_n+S\sin\Psi|<0.77S,
\]
where \(Y_j\) denotes the common positive normalization and
\(\Psi=(n+1)\theta-\Im G(w_1)\).
At an old zero, \(|\sin\Psi|\ge\sqrt{1-0.20^2}>0.77\).
The exterior join and the first hits of the integer phase levels
give the root count and the actual next-row signs.
\href{https://github.com/ProgramComputer/open-math-research/tree/main/topics/long-cycle-excedances/fourier-circle}{Proof and exact replay}

The two aspect intervals overlap on \([44,48]\). For \(n/k<4\)
the three-cycle theorem applies. Together these cover all admissible
thresholds at every old size from \(10{,}000\).
